\documentclass[10pt,a4paper]{amsart}
\usepackage{amsmath,amssymb,amsthm,graphicx}
\usepackage{verbatim}
\newenvironment{reference}{\comment}{\endcomment}
%\newenvironment{reference}{}{}
\newenvironment{slogan}{\comment}{\endcomment}
\newenvironment{history}{\comment}{\endcomment}

% For commutative diagrams we use Xy-pic
\usepackage[all]{xy}

% We use 2cell for 2-commutative diagrams.
\xyoption{2cell}
% Theorem environments.
%
\theoremstyle{plain}
\newtheorem{theorem}[subsection]{Theorem}
\newtheorem{proposition}[subsection]{Proposition}
\newtheorem{lemma}[subsection]{Lemma}

\theoremstyle{definition}
\newtheorem{definition}[subsection]{Definition}
\newtheorem{example}[subsection]{Example}
\newtheorem{exercise}[subsection]{Exercise}
\newtheorem{situation}[subsection]{Situation}

\theoremstyle{remark}
\newtheorem{remark}[subsection]{Remark}
\newtheorem{remarks}[subsection]{Remarks}

\numberwithin{equation}{subsection}

% Macros
%
\def\lim{\mathop{\mathrm{lim}}\nolimits}
\def\colim{\mathop{\mathrm{colim}}\nolimits}
\def\Spec{\mathop{\mathrm{Spec}}}
\def\Hom{\mathop{\mathrm{Hom}}\nolimits}
\def\Ext{\mathop{\mathrm{Ext}}\nolimits}
\def\SheafHom{\mathop{\mathcal{H}\!\mathit{om}}\nolimits}
\def\SheafExt{\mathop{\mathcal{E}\!\mathit{xt}}\nolimits}
\def\Sch{\mathit{Sch}}
\def\Mor{\mathop{\mathrm{Mor}}\nolimits}
\def\Ob{\mathop{\mathrm{Ob}}\nolimits}
\def\Sh{\mathop{\mathit{Sh}}\nolimits}
\def\NL{\mathop{N\!L}\nolimits}
\def\CH{\mathop{\mathrm{CH}}\nolimits}
\def\proetale{{pro\text{-}\acute{e}tale}}
\def\etale{{\acute{e}tale}}
\def\QCoh{\mathit{QCoh}}
\def\Ker{\mathop{\mathrm{Ker}}}
\def\Im{\mathop{\mathrm{Im}}}
\def\Coker{\mathop{\mathrm{Coker}}}
\def\Coim{\mathop{\mathrm{Coim}}}

% Boxtimes
%
\DeclareMathSymbol{\boxtimes}{\mathbin}{AMSa}{"02}

%
% Macros for moduli stacks/spaces
%
\def\QCohstack{\mathcal{QC}\!\mathit{oh}}
\def\Cohstack{\mathcal{C}\!\mathit{oh}}
\def\Spacesstack{\mathcal{S}\!\mathit{paces}}
\def\Quotfunctor{\mathrm{Quot}}
\def\Hilbfunctor{\mathrm{Hilb}}
\def\Curvesstack{\mathcal{C}\!\mathit{urves}}
\def\Polarizedstack{\mathcal{P}\!\mathit{olarized}}
\def\Complexesstack{\mathcal{C}\!\mathit{omplexes}}
% \Pic is the operator that assigns to X its picard group, usage \Pic(X)
% \Picardstack_{X/B} denotes the Picard stack of X over B
% \Picardfunctor_{X/B} denotes the Picard functor of X over B
\def\Pic{\mathop{\mathrm{Pic}}\nolimits}
\def\Picardstack{\mathcal{P}\!\mathit{ic}}
\def\Picardfunctor{\mathrm{Pic}}
\def\Deformationcategory{\mathcal{D}\!\mathit{ef}}
\usepackage{cleveref}
\crefname{originallemma}{original lemma}{original lemmas}
\Crefname{originallemma}{Original lemma}{Original lemmas}
\makeatletter
\newlabel{divisors-lemma-codim-1-part}{{0AGB}{source}{}{sourcepointer-1}{}}
\newlabel{divisors-lemma-codim-1-part@cref}{{[originallemma][0][]0AGB}{[1][0][]source}{}{}{}}
\makeatother
\crefname{originalreference}{original reference}{original references}
\Crefname{originalreference}{Original reference}{Original references}
\makeatletter
\newlabel{lemma-make-ideal-principal}{{0AHH}{source}{}{sourcepointer-2}{}}
\newlabel{lemma-make-ideal-principal@cref}{{[originalreference][0][]0AHH}{[1][0][]source}{}{}{}}
\makeatother
\title{Turning a proper closed subscheme of a regular surface into an effective Cartier divisor by blowups}
\author{The Stacks Project Authors}
\date{stacks-project:resolve.tex:local-snapshot-20261012; GNU FDL 1.2 or later}
\begin{document}
\maketitle
\noindent\textit{Editorial note.} All delivered mathematical text is the author's own, copied from the pinned Stacks Project chapter file resolve.tex; no statement or proof step was rewritten, reordered or removed.\par\medskip
\noindent\textit{Editorial note.} The counted proof is the author's complete argument: the decomposition of the ideal into the Cartier part and the residual part supported in finitely many points, and the principalization of the residual ideal.\par\medskip
\noindent\textit{Editorial note.} The label lemma-turn-into-effective-Cartier names the counted statement, and the proof environment immediately following it in the source is delivered verbatim as its proof.\par\medskip
\noindent\textit{Editorial note.} The counted statement is carried by the author's own lemma environment in the chapter source.\par\medskip
\noindent\textit{Editorial note.} Pre-open gate: python operations/agent-authoring/tools/precard\_gate.py --source resolve.tex --label lemma-turn-into-effective-Cartier --start 4423 --end 4445 --local-id 28118 returned RESULT=OK (L28118 free, occupancy FREE, public interval clear). Pre-assembly tag self-check: canonical tag resolved and not present in reference/public-items.jsonl.\par\medskip
\noindent\textit{Editorial note.} The statement and proof are copied from the pinned chapter file; the cited inputs are kept as external pointers with their canonical Stacks tags.\par\medskip
\noindent\textit{Allowed standard prerequisites (editorial).}\begin{itemize}
\item Decomposition of a closed subscheme into its largest Cartier divisor and a residual part of codimension two. Divisors, the lemma describing the codimension one part of a closed subscheme. Decomposes the ideal sheaf of the given subscheme as a product of the ideal of its largest effective Cartier divisor and the ideal of a subscheme supported in finitely many points.
\item Making an ideal defining finitely many points invertible by blowups in closed points. Same chapter, the lemma producing a sequence of blowups in closed points after which an ideal supported on finitely many regular surface points becomes invertible. Provides the sequence of blowups making the residual ideal invertible, which is what makes the inverse image of the whole subscheme an effective Cartier divisor.
\end{itemize}\medskip
\ifdefined\hypertarget\hypertarget{sourcepointer-1}{}\fi
\noindent\textit{Source reference.} Original lemma 0AGB: sources/stacks-chow/tags.txt, source lines 13189--13189. Prerequisite pointer invoked at proof line(s) 4437; it names a standard fact of a sibling chapter that the counted argument uses, and every step of the counted proof built on it is retained verbatim in the delivered body.\par\medskip
\ifdefined\hypertarget\hypertarget{sourcepointer-2}{}\fi
\noindent\textit{Source reference.} Original reference 0AHH: sources/stacks-resolve/pinned-original-source.tex, source lines 1062--1062. Prerequisite pointer invoked at proof line(s) 4441; it names a standard fact of the same chapter that the counted argument uses, and every step of the counted proof built on it is retained verbatim in the delivered body.\par\medskip
\begin{lemma}
\label{lemma-turn-into-effective-Cartier}
Let $X$ be a regular scheme of dimension $2$. Let $Z \subset X$
be a proper closed subscheme. There exists a sequence
$$
X_n \to \ldots \to X_1 \to X
$$
of blowing ups in closed points such that the inverse image $Z_n$ of $Z$
in $X_n$ is an effective Cartier divisor.
\end{lemma}

\begin{proof}
Let $D \subset Z$ be the largest effective Cartier divisor contained in $Z$.
Then $\mathcal{I}_Z \subset \mathcal{I}_D$ and the quotient is supported
in closed points by Divisors, Lemma \ref{divisors-lemma-codim-1-part}.
Thus we can write $\mathcal{I}_Z = \mathcal{I}_{Z'} \mathcal{I}_D$
where $Z' \subset X$ is a closed subscheme which set theoretically
consists of finitely many closed points. Applying
Lemma \ref{lemma-make-ideal-principal}
we find a sequence of blowups as in the statement of our lemma
such that $\mathcal{I}_{Z'}\mathcal{O}_{X_n}$ is invertible.
This proves the lemma.
\end{proof}

\end{document}
